\documentclass[a4paper, noanswers, addpoints, 12pt]{exam}
\usepackage[T1]{fontenc}
\usepackage[utf8]{inputenc}
\usepackage{lmodern}
\usepackage[french]{babel}
\usepackage{xcolor}
\usepackage{listings}
\usepackage{graphicx}
\usepackage{array}
\usepackage{eurosym}

% Configuration de la coloration syntaxique (style du document joint)
\definecolor{codegreen}{rgb}{0,0.6,0}
\definecolor{codegray}{rgb}{0.5,0.5,0.5}
\definecolor{codepurple}{rgb}{0.58,0,0.82}
\definecolor{backcolour}{rgb}{0.95,0.95,0.92}
 
\lstdefinestyle{mystyle}{
    backgroundcolor=\color{backcolour},   
    commentstyle=\color{codegreen},
    keywordstyle=\color{magenta},
    numberstyle=\tiny\color{codegray},
    stringstyle=\color{codepurple},
    basicstyle=\footnotesize\ttfamily,
    breakatwhitespace=false,         
    breaklines=true,                 
    captionpos=b,                    
    keepspaces=true,                 
    numbers=left,                    
    numbersep=5pt,                  
    showspaces=false,                
    showstringspaces=false,
    showtabs=false,                  
    tabsize=2
}
\lstset{style=mystyle, language=Python}

\title{\vspace{-2cm}Examen Première NSI - Types construits et gloutonnerie}
\date{}

\usepackage[a4paper,top=2.3cm,bottom=2.8cm,left=2cm,right=2cm,headsep=0.6cm,footskip=1.2cm]{geometry}
\begin{document}
\runningfooter{}{\footnotesize Babash - NSI Tani Malandi et Johanne Bakalara \textperiodcentered{} CC BY-SA 4.0 \textperiodcentered{} page \thepage}{}
\firstpagefooter{}{\footnotesize Babash - NSI Tani Malandi et Johanne Bakalara \textperiodcentered{} CC BY-SA 4.0 \textperiodcentered{} page \thepage}{}

\maketitle
\thispagestyle{headandfoot}

\pagestyle{headandfoot}
\thispagestyle{headandfoot}
\firstpageheader{\makebox[\textwidth]{Nom : \dotfill\ Pr\'enom : \dotfill\ Classe : \makebox[3cm]{\dotfill}}}{}{}
\vspace{-3cm}


\section*{Consignes:}
\begin{center} 
\fbox{\parbox{\linewidth}{
\begin{center}
La clarté de la rédaction et les explications font partie intégrante de la notation. Pour les questions demandant du code, respectez la syntaxe Python et soyez attentifs aux détails.
\end{center}


}}
\end{center}

\vsword{Note}
\hsword{Note}

\section*{Barème :}
\begin{center}
  \par\bigskip
  \gradetable[h]
\end{center}

\vspace{0.5cm}

\begin{questions}

% --- EXERCICE 1 : STRUCTURES ET DICTIONNAIRES (9,5 points) ---
\question \textbf{Maîtrise des types construits}

\begin{parts}
    \part[2] Pour chaque besoin, indiquez la structure (\textbf{Liste} ou \textbf{Dictionnaire}) la plus adaptée. Justifiez (une réponse sans justification sera considérée fausse \textit{0,5 pt par réponse}).
    \begin{itemize}
        \item Conserver l'ordre d'arrivée de clients dans une file d'attente : \dotfill\\
        Justification : \dotfill \\
        .\dotfill
        \item Associer chaque département français (ex: "33") à son nom ("Gironde") : \dotfill\\
        Justification : \dotfill \\
        .\dotfill
        \item Stocker les températures relevées toutes les heures dans une journée : \dotfill\\
        Justification : \dotfill \\
        .\dotfill
        \item Créer un inventaire de jeu vidéo pour chercher un objet par son nom : \dotfill\\
        Justification : \dotfill \\
        .\dotfill
    \end{itemize}

    \part[1] \textbf{Création :} Écrivez une instruction Python pour créer un dictionnaire nommé \texttt{mon\_perso} possédant deux clés : \texttt{"nom"} (valeur "Arthur") et \texttt{"vie"} (valeur 100).\\
    .\dotfill\\
    .\dotfill

    \part[1] \textbf{Accès :} Soit \texttt{livres = \{"Dune": 3, "1984": 2\}}. Écrivez l'instruction pour récupérer le stock de "1984" dans une variable \texttt{n}.\\
    .\dotfill

    \part[1] \textbf{Modification :} Écrivez l'instruction pour modifier le stock de "Dune" à 4 dans le dictionnaire précédent.\\
    .\dotfill

    \part[1] \textbf{Recherche :} Écrivez une instruction permettant de tester si la clé \texttt{"Zelda"} est présente dans \texttt{livres} (doit renvoyer un booléen).\\
    .\dotfill

    \part[\half] \textbf{Analyse de boucle :} Soit le code suivant :
        \begin{lstlisting}[language=Python]
    for k in livres.keys():
        print(k)\end{lstlisting}
        Quelle catégorie d'informations (les titres, les quantités ou les deux) sera affichée par ce programme ?\dotfill
    
    \part[1] \textbf{Analyse de boucle :} On souhaite calculer le nombre total de livres. Complétez la ligne manquante :
    \begin{lstlisting}[language=Python]
total = 0
for v in ................................... :
    total = total + v\end{lstlisting}

    \part[\half] \textbf{Analyse de boucle :} Si on utilise la boucle \texttt{for a, b in livres.items():}, que contiendra la variable \texttt{a} à chaque tour de boucle ?
    \begin{itemize}
        \item tour 0 :\dotfill
        \item tour 1 :\dotfill
    \end{itemize}

   \part[2] \textbf{Parcours :} Écrivez un script (boucle \texttt{for}) qui parcourt le dictionnaire \texttt{livres} et affiche le titre de chaque livre dont le stock est $>= 2$.\\
    .\dotfill\\
    .\dotfill\\
    .\dotfill\\
    .\dotfill\\
    .\dotfill

\end{parts}

\vspace{0.4cm}

% --- EXERCICE 2 : STRUCTURES COMPOSITES (3 points) ---
\question \textbf{Analyse d'un tournoi d'E-sport}
\begin{lstlisting}
scores = [
    {"pseudo": "Alpha", "points": 1500, "serveur": "EUW"},
    {"pseudo": "Beta", "points": 2800, "serveur": "NA"},
    {"pseudo": "Gamma", "points": 1200, "serveur": "EUW"}
]
\end{lstlisting}

\begin{parts}
    \part[1] Quelle instruction permet d'accéder au \textbf{pseudo} du deuxième joueur de la liste (Beta) ?\\
    .\dotfill
    
    \part[1] \textbf{Utilisation de len() :}\\ Que renvoie l'instruction : \texttt{len(scores[0]["pseudo"])}?\dotfill\\
    Justification : \dotfill \\
    
    \part[2] Complétez la fonction ci-dessous pour qu'elle renvoie la somme des points des joueurs du serveur $zone$.
\begin{lstlisting}
def total_points_zone(liste_joueurs, zone):
    total = 0
    for joueur in ......................... :
        if joueur["serveur"] == zone:
            total = total + .........................
    return total
\end{lstlisting}
\end{parts}

% --- EXERCICE 3 : ALGORITHMES GLOUTONS (7,5 points) ---
\question \textbf{Optimisation de chargement (Sac à dos)}

Un transporteur doit charger son camion (capacité maximale : \textbf{120 kg}).
\begin{center}
\begin{tabular}{|l|c|c|c|c|}
\hline
Colis & A & B & C & D \\ \hline
Poids (kg) & 80 & 60 & 40 & 20 \\ \hline
Valeur (\euro) & 180 & 110 & 80 & 90 \\ \hline
\end{tabular}
\end{center}

\begin{parts}
    \part[1] \textbf{Analyse d'erreur :} Un élève tape l'instruction \texttt{print(colis["a"])} pour accéder au premier colis du tableau.\\
    Expliquez pourquoi Python renvoie une erreur de type \texttt{KeyError}\\
    .\dotfill\\
    .\dotfill\\

    \part[1] \textbf{Stratégie 1 :} En choisissant le "plus cher d'abord" (\textit{Capacité : 120 kg}), quels colis sont embarqués et quelle est la valeur totale ?\\
    .\dotfill\\
    .\dotfill\\

    \part[1] \textbf{Stratégie 2 :} Calculez les ratios (Valeur/Poids) des colis B et C. Lequel de ces deux colis la stratégie "Meilleur Ratio" choisirait-elle en priorité ?\\
    .\dotfill\\
    .\dotfill\\
    .\dotfill

    \part[1] Proposez une combinaison de colis respectant les \textbf{120 kg} dont la valeur est supérieure à celle obtenue par la stratégie 1. Justifiez l'intérêt.\\
    .\dotfill\\
    .\dotfill\\

    \part[2] \textbf{Synthèse :} Complétez le texte avec : \textit{garantit, ne garantit pas, couteuse, rapide, local(e), optimal(e), global(e)}.\\
    
    
    L'algorithme glouton est une méthode très ......................... car il ne revient jamais en arrière. À chaque étape, il fait le meilleur choix ......................... sans prévoir la suite. Cependant, cet algorithme ................................................................ toujours de trouver la solution ......................... ......................... .
\end{parts}

\end{questions}

\begin{center}
    \rule{0.5\textwidth}{0.4pt}\\
    \footnotesize Fin de l'énoncé.
\end{center}

\end{document}